\section*{Capítulo \ref{ch:g12:continental-crust} --- La corteza continental}

\begin{solution}{exo:g12:continental-crust:1}
\omterm{def:g12:continental-crust:continental-crust}{Corteza continental}: de 30 a \qty{70}{km}, unos \qty{41}{km} de media; rocas
graníticas, gneises y otras rocas metamórficas y \omterm{def:g11:igneous-rocks:volcanic-plutonic}{plutónicas} bajo una
cobertera sedimentaria; un \qty{61}{\percent} de sílice de media; unos
\qty{2830}{kg/m^3}; hasta 4030 millones de años (granos de 4400). \omterm{def:g11:ocean-ridges:oceanic-lithosphere}{Corteza
oceánica}: unos \qty{7}{km}; basalto y \omterm{def:g11:ocean-ridges:gabbro}{gabro}; un \qty{50}{\percent} de sílice;
unos \qty{3000}{kg/m^3}; menos de 200 millones de años.
\end{solution}

\begin{solution}{exo:g12:continental-crust:2}
\qty{60.6}{\percent}: \omterm{def:g11:igneous-rocks:felsic-mafic}{intermedia}, como una \omterm{def:g11:igneous-rocks:rock-names}{diorita}. \qty{66.6}{\percent}:
\omterm{def:g11:igneous-rocks:felsic-mafic}{félsica}, como un granito.
\end{solution}

\begin{solution}{exo:g12:continental-crust:3}
La \omterm{def:g11:ocean-ridges:oceanic-lithosphere}{corteza oceánica}, de unos \qty{3000}{kg/m^3}, descansa sobre manto
litosférico, y la \omterm{def:g11:ocean-ridges:oceanic-lithosphere}{litosfera oceánica} vieja y fría es, en conjunto, lo bastante
densa para hundirse. La \omterm{def:g12:continental-crust:continental-crust}{corteza continental}, de unos \qty{2830}{kg/m^3},
mucho más ligera que el manto (unos \qty{3380}{kg/m^3}), flota demasiado: un
continente que llega a una fosa detiene la \omterm{def:g11:subduction-zones:subduction}{subducción}.
\end{solution}

\begin{solution}{exo:g12:continental-crust:4}
(1) Un núcleo antiguo y estable de un continente, que ya no se deforma. (2)
Más de mil millones de años. (3) Rocas metamórficas y \omterm{def:g11:igneous-rocks:volcanic-plutonic}{plutónicas}, las raíces
profundas de antiguas montañas.
\end{solution}

\begin{solution}{exo:g12:continental-crust:5}
$4030/190 \approx 21$ veces. El fondo oceánico se forma en las dorsales y se
destruye en las \omterm{def:g11:subduction-zones:subduction}{zonas de subducción}: nada de él sobrevive mucho más de 200
millones de años.
\end{solution}

\begin{solution}{exo:g12:continental-crust:6}
(1) $0.306 \times 210.4 \approx \num{64.4}\text{ millones de }\unit{km^2}$.
(2) $210.4 - 64.4 = \num{146.0}\text{ millones de }\unit{km^2}$.
(3) $4\pi \times 6371^2 \approx \qty{5.10e8}{km^2}$; $210.4/510.1 \approx
0.41$: el \qty{41}{\percent}.
\end{solution}

\begin{solution}{exo:g12:continental-crust:7}
(1) $41100 \times 2830 \approx \qty{1.16e8}{kg}$. (2) $7100 \times 3000 =
\qty{2.13e7}{kg}$. (3) Unas 5.5 veces.
\end{solution}

\begin{solution}{exo:g12:continental-crust:8}
\omterm{def:g12:continental-crust:root}{Raíz}: $2.5 \times 2830/550 \approx \qty{12.9}{km}$. La corteza es más gruesa en
$2.5 \times 3380/550 \approx \qty{15.4}{km}$: unos \qty{50}{km} en total,
dentro de los 50 a \qty{70}{km} medidos bajo los cinturones montañosos.
\end{solution}

\begin{solution}{exo:g12:continental-crust:9}
(1) Alrededor de $0.103 + 0.100 + 0.102 \approx 0.30$: el \qty{30}{\percent}.
(2) La banda más joven, de 0 a 10 millones de años (alrededor del
\qty{10}{\percent}), casi empatada con la de 20 a 30. (3) El fondo antiguo se
ha destruido en las \omterm{def:g11:subduction-zones:subduction}{zonas de subducción}; cuanto más antiguo es el fondo, más
parte de él ha desaparecido.
\end{solution}

\begin{solution}{exo:g12:continental-crust:10}
$120 = 0.694\,h - 0.306 \times 100$, así que $h = 150.6/0.694 \approx
\qty{217}{m}$. Contar solo las tierras emergidas deja fuera los bordes bajos y
sumergidos de los continentes, y la media sale mucho más alta.
\end{solution}

\begin{solution}{exo:g12:continental-crust:11}
La isla: un \omterm{def:g11:ocean-ridges:gabbro}{gabro} bajo lavas y un \omterm{def:g11:seismic-structure:moho}{Moho} a \qty{8}{km}, \omterm{def:g11:ocean-ridges:oceanic-lithosphere}{corteza oceánica}
(la altitud de la isla no cambia eso). La llanura: gneis y un
\omterm{def:g11:seismic-structure:moho}{Moho} a \qty{38}{km}, \omterm{def:g12:continental-crust:continental-crust}{corteza continental}.
\end{solution}

\begin{solution}{exo:g12:continental-crust:12}
(1) Exacto: $\tfrac{4}{3}\pi(6371^3 - 6329.9^3) \approx
\qty{2.0829e10}{km^3}$; aproximado: $4\pi \times 6371^2 \times 41.1
\approx \qty{2.0964e10}{km^3}$. (2) La aproximación es demasiado grande en un
\qty{0.65}{\percent}, cerca de $H/R \approx 0.0065$: para una capa delgada el
error es despreciable.
\end{solution}

\begin{solution}{exo:g12:continental-crust:13}
(1) $60.6 = (66.6 + 2x)/3$, así que $x = (181.8 - 66.6)/2 = \qty{57.6}{\percent}$.
(2) \omterm{def:g11:igneous-rocks:felsic-mafic}{Intermedia}. (3) La \omterm{def:g12:continental-crust:continental-crust}{corteza continental} profunda es menos rica en sílice
que la corteza superior: más minerales \omterm{def:g11:igneous-rocks:felsic-mafic}{máficos}, más rica en hierro, magnesio y
calcio.
\end{solution}

\begin{solution}{exo:g12:continental-crust:14}
(1) $3 \times 10^8/3 = 10^8$ años, 100 millones de años. (2) Del mismo orden
que las edades del fondo oceánico: tres cuartas partes tienen menos de 100
millones de años, y ninguna más de unos 190. (3) Tiene que destruirse tan
deprisa como se forma: se hunde de nuevo en el manto en las \omterm{def:g11:subduction-zones:subduction}{zonas de
subducción}.
\end{solution}

\begin{solution}{exo:g12:continental-crust:15}
(1) $8.6 \times 10^9 / 4.0 \times 10^9 \approx \qty{2.2}{km^3}$ al año.
(2) Un ritmo constante es solo una posibilidad: la corteza pudo crecer deprisa
al principio y reciclarse después, o crecer tarde. Las rocas y los granos más
antiguos prueban que pronto existía algo de corteza; no pueden decir cuánta,
porque la corteza antigua puede erosionarse o volver al manto y dejar pocas
huellas.
\end{solution}

\begin{solution}{pb:g12:continental-crust:1}
\textbf{1.} $4\pi \times 6371^2 \approx \qty{5.10e8}{km^2}$.

\textbf{2.} $210.4/510.1 \approx 0.41$: el \qty{41}{\percent}.

\textbf{3.} $0.306 \times 210.4 \approx \num{64.4}\text{ millones de }\unit{km^2}$.

\textbf{4.} $510.1 - 210.4 \approx \num{299.7}\text{ millones de }\unit{km^2}$,
el \qty{59}{\percent}.

\textbf{5.} Los bordes de los continentes están bajo mares poco profundos, en
las plataformas continentales.

\textbf{6.} El espesor, \qty{41}{km}, es muy pequeño comparado con el radio,
\qty{6371}{km}: la curvatura cambia el volumen en menos de un uno por
ciento.

\textbf{7.} $210.4 \times 10^6 \times 41.1 \approx \qty{8.65e9}{km^3}$.

\textbf{8.} \qty{8.65e18}{m^3}.

\textbf{9.} $210.4 \times 10^6 \times 36 \approx \qty{7.57e9}{km^3}$.

\textbf{10.} $(8.65 - 7.57)/7.57 \approx 0.14$: \qty{14}{\percent}.

\textbf{11.} $8.65 \times 10^{18} \times 2830 \approx \qty{2.45e22}{kg}$.

\textbf{12.} $41100 \times 2830 \approx \qty{1.16e8}{kg}$.

\textbf{13.} $299.7 \times 10^6 \times 7.1 \approx \qty{2.13e9}{km^3}$, es
decir, \qty{2.13e18}{m^3}; por 3000: unos \qty{6.4e21}{kg}.

\textbf{14.} La \omterm{def:g12:continental-crust:continental-crust}{corteza continental}, unas 3.8 veces más pesada: es casi seis
veces más gruesa, lo que compensa con creces su menor superficie y su menor
densidad.

\textbf{15.} Está formada por rocas graníticas, ricas en cuarzo y feldespatos,
los minerales formadores de rocas más ligeros; la \omterm{def:g11:ocean-ridges:oceanic-lithosphere}{corteza oceánica} está formada
por basalto y \omterm{def:g11:ocean-ridges:gabbro}{gabro}, ricos en piroxeno y olivino.

\textbf{16.} $2.45 \times 10^{22}/5.97 \times 10^{24} \approx 0.0041$.

\textbf{17.} $3 \times 0.41 \times (41.1/6371) \times (2830/5513) \approx
0.0041$: lo mismo.

\textbf{18.} $6.4 \times 10^{21}/5.97 \times 10^{24} \approx 0.0011$: alrededor
de una décima de por ciento.

\textbf{19.} Alrededor del \qty{0.4}{\percent}.

\textbf{20.} La \omterm{def:g12:continental-crust:continental-crust}{corteza continental} es alrededor del \textbf{\qty{0.4}{\percent} de la
masa de la Tierra}.
\end{solution}
